% ============================================================
%  MATH 5345H Fall 2026 — Combined Lecture Notes
%  Replace this file with your actual lecture notes.
%
%  WEEK BOUNDARY MARKERS (required by the build script):
%    %%%WEEK N | Short title
%  Place this comment on its own line immediately before the
%  first content of week N.  Everything above \begin{document}
%  is the shared preamble and is prepended to every weekly file.
% ============================================================

\documentclass{article}
\usepackage{amsmath, amsthm, amssymb}
\usepackage{hyperref}

\newtheorem{theorem}{Theorem}
\newtheorem{lemma}[theorem]{Lemma}
\newtheorem{proposition}[theorem]{Proposition}
\newtheorem{corollary}[theorem]{Corollary}
\newtheorem{definition}[theorem]{Definition}
\newtheorem{example}[theorem]{Example}
\newtheorem{remark}[theorem]{Remark}

\title{MATH 5345H: Honors Introduction to Topology\\Fall 2026}
\author{Erkao Bao}
\date{}

\begin{document}
\maketitle

%%%WEEK 1 | Course overview; sets, functions, relations

\section*{Week 1 — Course Overview and Foundations}

% Replace the placeholder text below with your actual Week 1 notes.

\begin{definition}
A \emph{set} is a collection of distinct objects, called \emph{elements}.
\end{definition}

%%%WEEK 2 | Cartesian products; countable and uncountable sets

\section*{Week 2 — Cartesian Products and Cardinality}

% Replace with actual Week 2 notes.

\begin{definition}
Given sets $A$ and $B$, the \emph{Cartesian product} $A \times B$ is the set of all ordered pairs $(a,b)$ with $a \in A$ and $b \in B$.
\end{definition}

%%%WEEK 3 | Topological spaces; bases and subbases

\section*{Week 3 — Topological Spaces}

% Replace with actual Week 3 notes.

\begin{definition}
A \emph{topological space} is a set $X$ together with a collection $\mathcal{T}$ of subsets of $X$, called \emph{open sets}, satisfying:
\begin{enumerate}
  \item $\emptyset \in \mathcal{T}$ and $X \in \mathcal{T}$.
  \item Arbitrary unions of members of $\mathcal{T}$ belong to $\mathcal{T}$.
  \item Finite intersections of members of $\mathcal{T}$ belong to $\mathcal{T}$.
\end{enumerate}
\end{definition}

%%%WEEK 4 | Order topology; product topology; subspaces

\section*{Week 4 — Order and Product Topologies}

% Replace with actual Week 4 notes.

%%%WEEK 5 | Continuity; homeomorphisms; arbitrary products

\section*{Week 5 — Continuity and Homeomorphisms}

% Replace with actual Week 5 notes.

%%%WEEK 6 | Metric topologies; metrizability; sequences

\section*{Week 6 — Metric Topologies}

% Replace with actual Week 6 notes.

%%%WEEK 7 | Quotient topology; connectedness

\section*{Week 7 — Quotient Topology and Connectedness}

% Replace with actual Week 7 notes.

%%%WEEK 8 | Path components; compactness

\section*{Week 8 — Compactness}

% Replace with actual Week 8 notes.

%%%WEEK 9 | Heine-Borel; limit point and sequential compactness

\section*{Week 9 — Compactness in Euclidean Space}

% Replace with actual Week 9 notes.

%%%WEEK 10 | Local compactness; countability and separation axioms

\section*{Week 10 — Local Compactness and Separation Axioms}

% Replace with actual Week 10 notes.

%%%WEEK 11 | Urysohn lemma; metrization; Tietze extension

\section*{Week 11 — Urysohn Lemma and Metrization}

% Replace with actual Week 11 notes.

%%%WEEK 12 | Tychonoff theorem; Nagata-Smirnov

\section*{Week 12 — The Tychonoff Theorem}

% Replace with actual Week 12 notes.

%%%WEEK 13 | Paracompactness; Smirnov; complete metric spaces

\section*{Week 13 — Paracompactness and Complete Metric Spaces}

% Replace with actual Week 13 notes.

%%%WEEK 14 | Compact convergence; Ascoli; Baire spaces

\section*{Week 14 — Function Spaces and Baire Spaces}

% Replace with actual Week 14 notes.

%%%WEEK 15 | Homotopy of paths; course review

\section*{Week 15 — Homotopy of Paths}

% Replace with actual Week 15 notes.

\end{document}
